\subsection{复形态射的锥与柱}

设 \(u:(\mathbf{C}^{\prime},d^{\prime})\to(\mathbf{C},d)\) 是复形的一个态射。设 \(\operatorname{Cyl}(u)\) 与 \(\operatorname{Con}(u)\) 为 \(\text{分次 A-模 }\operatorname{Cyl}(u)=\mathbf{C}^{\prime}\oplus\mathbf{C}^{\prime}(-1)\oplus\mathbf{C}\) , \(\operatorname{Con}(u)=\mathbf{C}^{\prime}(-1)\oplus\mathbf{C}\) ，并让我们定义次数为 \((-1)\) 的分次 A-线性映射

\(\overline{\mathbf{D}}:\operatorname {Cyl}(u)\to \operatorname {Cyl}(u),\quad \overline{\mathbf{D}} (x^{\prime},y^{\prime},x) = (d^{\prime}x^{\prime} + y^{\prime}, - d^{\prime}y^{\prime},dx - u(y^{\prime})),\)

D : Con \((u) \to \text{Con}(u)\) , D \((y', x) = (-d'y', dx - u(y'))\) .

（此处以及后文中，我们用 \(x\) , \(y\) , \ldots{} 表示 C 的任意元素，用 \(x'\) , \(y'\) , \ldots{} 表示 C' 的任意元素。）

引理2. --- (Cyl \((u)\) , \(\overline{\mathbf{D}}\) ) 和 (Con \((u)\) , D) 是A-模的复形。

确实，我们有

\[
\begin{array}{r l} \overline {{\mathbf {D}}} \circ \overline {{\mathbf {D}}} (x ^ {\prime}, y ^ {\prime}, x) & = \overline {{\mathbf {D}}} (d ^ {\prime} x ^ {\prime} + y ^ {\prime}, - d ^ {\prime} y ^ {\prime}, d x - u (y ^ {\prime})) = \\ & = (d ^ {\prime} (d ^ {\prime} x ^ {\prime} + y ^ {\prime}) - d ^ {\prime} y ^ {\prime}, - d ^ {\prime} (- d ^ {\prime} y ^ {\prime}), d (d x - u (y ^ {\prime})) - u (- d ^ {\prime} y ^ {\prime})) = 0 \end{array}
\]

因为 \(d' \circ d' = 0\) , \(d \circ d = 0\) 且 \(d \circ u = u \circ d'\) 。同样地 \(D \circ D = 0\) .

定义7.——复形Cyl(u)和Con(u)分别称为态射u的柱和锥。

例. --- 设 \(u: \mathbf{M} \to \mathbf{N}\) 是 A-模的同态；则 \(\operatorname{Cyl}(u)\) 与 \(\operatorname{Con}(u)\) 的仅有的非零齐次分量是

\[
\operatorname{Cyl} _ {1} (u) = \mathrm{M}, \quad \operatorname{Cyl} _ {0} (u) = \mathrm{M} \oplus \mathrm{N},
\]

\[
\operatorname{Con} _ {1} (u) = \mathrm{M}, \quad \operatorname{Con} _ {0} (u) = \mathrm{N},
\]

并且我们有 \(\overline{\mathrm{D}}(m) = (m, -u(m))\) , \(\mathrm{D}(m) = -u(m)\) （对 \(m \in \mathbf{M}\) ）；因此，

\[
\mathrm{H} _ {1} (\text { Con } (u)) = \operatorname{Ker} (u), \quad \mathrm{H} _ {0} (\text { Con } (u)) = \operatorname{Coker} (u).
\]

复形 \(\Gamma(s)\) （与空间 \(\{s\}\) 相关联，后者赋予其唯一的胞腔分解）约化为模 A，并被等同于 \(\Gamma(\text{Con}(f))\) 的一个子复形；用 \(\Gamma(\text{Con}(f), s)\) 表示商复形。映射 \(f\) 定义了复形的一个态射 \(\Gamma(f): \Gamma(X) \to \Gamma(Y)\) ，并且可以证明，复形 \(\Gamma(\text{Cyl}(f))\) 与 \(\Gamma(\text{Con}(f), s)\) 分别被等同于 \(\text{Cyl}(\Gamma(f))\) 与 \(\text{Con}(\Gamma(f))\) .

此外，在不对 X 和 Y 作任何假设的情况下，可将 \(f\) 与复形的一个态射 \(f_{*}: \mathrm{C}(\mathrm{X}, \mathrm{A}) \to \mathrm{C}(\mathrm{Y}, \mathrm{A})\) 相关联。可以构造从 \(\mathrm{Cyl}(f_{*})\) 到 \(\mathrm{C}(\mathrm{Cyl}(f), \mathrm{A})\) 以及从 \(\mathrm{Con}(f_{*})\) 到 \(\mathrm{C}(\mathrm{Con}(f), \{s\}, \mathrm{A})\) 的单射同伦等价。

设 \(\tilde{f}: X \to \text{Cyl}(f)\) 是把 \(x\) 映到 \((0, x)\) 的像的映射， \(\alpha: Y \to \text{Cyl}(f)\) 为典范映射， \(\beta: \text{Cyl}(f) \to Y\) 为这样的映射：它把 \(y\) 映到其在 \(\text{Cyl}(f)\) 中的像（当 \(y \in Y\) ），并把 \(f(x)\) 映到 \((t, x)\) 在 \(\text{Cyl}(f)\) 中的像（当 \(t \in [0, 1]\) 且 \(x \in X\) ）。映射 \(\tilde{f}\) 是从 \(X\) 到 \(\text{Cyl}(f)\) 的一个闭子集上的同胚，我们有 \(\beta \circ \alpha = \text{Id}_Y\) ，且 \(\alpha \circ \beta\) 在拓扑上同伦于 \(\text{Cyl}(f)\) 的恒等映射。这些性质应与下面的命题7进行比较。

现在考虑次数为 0 的分次 A-线性映射。

\[
\tilde {u}: \mathrm{C} ^ {\prime} \rightarrow \operatorname{Cyl} (u), \quad \tilde {u} (x ^ {\prime}) = (x ^ {\prime}, 0, 0),
\]

\[
\alpha : \mathrm{C} \rightarrow \operatorname{Cyl} (u), \quad \alpha (x) = (0, 0, x),
\]

\[
\beta : \operatorname{Cyl} (u) \rightarrow \mathrm{C}, \quad \beta \left(x ^ {\prime}, y ^ {\prime}, x\right) = u \left(x ^ {\prime}\right) + x,
\]

\[
\pi : \mathbf {C} \rightarrow \operatorname{Con} (u), \quad \pi (x) = (0, x),
\]

\[
\tilde {\pi}: \operatorname{Cyl} (u) \to \operatorname{Con} (u), \quad \tilde {\pi} (x ^ {\prime}, y ^ {\prime}, x) = (y ^ {\prime}, x),
\]

\[
\delta : \operatorname{Con} (u) \to \mathbf {C} ^ {\prime} (- 1), \qquad \delta (y ^ {\prime}, x) = y ^ {\prime}.
\]

命题7. --- a) 映射 \(\tilde{u}, \alpha, \beta, \pi, \tilde{\pi}, \delta\) 是复形的态射：我们有 \(u = \beta \circ \tilde{u}, \pi = \tilde{\pi} \circ \alpha, \beta \circ \alpha = 1_{\mathbb{C}}\) .

\begin{enumerate}
\def\labelenumi{\alph{enumi})}
\setcounter{enumi}{1}
\tightlist
\item
  复形态射的序列
\end{enumerate}

\begin{enumerate}
\def\labelenumi{(\arabic{enumi})}
\setcounter{enumi}{5}
\tightlist
\item
\end{enumerate}

\[
0 \rightarrow \mathrm{C} ^ {\prime} \stackrel {\tilde {u}} {\rightarrow} \operatorname{Cyl} (u) \stackrel {\tilde {\pi}} {\rightarrow} \operatorname{Con} (u) \rightarrow 0\tag{7}
\]

\[
0 \rightarrow \mathrm{C} \xrightarrow {\pi} \operatorname{Con} (u) \xrightarrow {\delta} \mathrm{C} ^ {\prime} (- 1) \rightarrow 0
\]

是正合的。

\begin{enumerate}
\def\labelenumi{\alph{enumi})}
\setcounter{enumi}{2}
\tightlist
\item
  态射 \(\alpha : \mathbf{C} \to \operatorname{Cyl}(u)\) 与 \(\beta : \operatorname{Cyl}(u) \to \mathbf{C}\) 是在同伦意义下互逆的同伦等价。
\end{enumerate}

断言a)等价于以下公式

\[
\tilde {u} \circ d ^ {\prime} = \overline {{\mathbf {D}}} \circ \tilde {u}, \alpha \circ d = \overline {{\mathbf {D}}} \circ \alpha , \beta \circ \overline {{\mathbf {D}}} = d \circ \beta , \pi \circ d = \mathbf {D} \circ \pi ,
\]

\[
\tilde {\pi} \circ \overline {{\mathbf {D}}} = \mathbf {D} \circ \tilde {\pi}, \delta \circ \mathbf {D} = - d ^ {\prime} \circ \delta , u = \beta \circ \tilde {u}, \pi = \tilde {\pi} \circ \alpha , \beta \circ \alpha = 1 _ {\mathrm{c}}
\]

这些通过直接计算即可验证。断言b)是平凡的。我们来证明c)；一方面有 \(\beta \circ \alpha = 1_{c}\) ；另一方面，若 \(\sigma: \text{Cyl}(u) \to \text{Cyl}(u)\) 是次数为1的A-线性分次映射，使得 \(\sigma(x', y', x) = (0, x', 0)\) ，人们立即验证

\[
\overline {{{D}}} \circ \sigma + \sigma \circ \overline {{{D}}} + \alpha \circ \beta = 1 _ {\mathrm{Cyl} (u)},
\]

由此得 c)。

可以用以下交换图来总结命题7，其中各行是正合的，各垂直箭头是同伦等价：

\[
\begin{array}{c c c c c} 0 & \longrightarrow & C & \stackrel {{\pi}} {{\longrightarrow}} & \operatorname{Con} (u) \stackrel {{\delta}} {{\longrightarrow}} C ^ {\prime} (- 1) \longrightarrow 0 \\ & & \Big \downarrow^ {\alpha} & & \Big \downarrow^ {1} \\ 0 & \longrightarrow & C ^ {\prime} \stackrel {{\bar {u}}} {{\longrightarrow}} & \operatorname{Cyl} (u) \stackrel {{\tilde {\pi}}} {{\longrightarrow}} & \operatorname{Con} (u) \longrightarrow 0 \\ & & \Big \downarrow^ {1} & & \Big \downarrow^ {\beta} \\ & & C ^ {\prime} \stackrel {{u}} {{\longrightarrow}} & C \end{array}
\]

推论。—— 对于每个复形态射 \(u: \mathbf{C}' \to \mathbf{C}\) ，存在一个单射的复形态射 \(\tilde{u}: \mathbf{C}' \to \mathbf{C}_1\) 以及一个同伦等价 \(\beta: \mathbf{C}_1 \to \mathbf{C}\) 使得 \(u = \beta \circ \tilde{u}\) .

引理3。—— a) 连接同态

\[
\partial_ {n + 1} (\pi , \delta): \mathrm{H} _ {n} (\mathrm{C} ^ {\prime}) \rightarrow \mathrm{H} _ {n} (\mathrm{C})
\]

相对于正合列(7)等于 -H \(_{n}\) (u)。

\begin{enumerate}
\def\labelenumi{\alph{enumi})}
\setcounter{enumi}{1}
\tightlist
\item
  连接同态
\end{enumerate}

\[
\partial_ {n} (\tilde {u}, \tilde {\pi}): \mathrm{H} _ {n} (\operatorname{Con} (u)) \rightarrow \mathrm{H} _ {n - 1} (\mathrm{C} ^ {\prime})
\]

相对于正合列(6)等于H \(_{n}\) (δ).

设 \(x' \in Z_n(\mathbf{C}')\) ；由于 \(x' = \delta(x', 0)\) 且

\[
- \mathrm{D} (x ^ {\prime}, 0) = (d ^ {\prime} x ^ {\prime}, u (x ^ {\prime})) = (0, u (x ^ {\prime})) = \pi (u (x ^ {\prime})),
\]

\(\partial (\pi ,\delta)\) 根据定义把 \(x^{\prime}\) 在 \(\mathbf{H}_n(\mathbf{C}')\) 中的类映到 \(-u(x^{\prime})\) 在 \(\mathbf{H}_n(\mathbf{C})\) 中的类，由此得 a)。

设 \((y', x) \in \operatorname{Con}_n(u)\) 使得 \(\mathbf{D}(y', x) = 0\) ；于是我们有 \((-d'y', dx - u(y')) = 0\) 。由于 \((y', x) = \tilde{\pi}(0, y', x)\) 且

\[
\overline {{{\mathrm{D}}}} (0, y ^ {\prime}, x) = \left(y ^ {\prime}, - d ^ {\prime} y ^ {\prime}, d x - u (y ^ {\prime})\right) = \left(y ^ {\prime}, 0, 0\right) = \tilde {u} \big (\delta (y ^ {\prime}, x) \big),
\]

\(\partial (\tilde{u},\tilde{\pi})\) 根据定义把 \((y^{\prime},x)\) 在 \(\mathrm{H}_n(\mathrm{Con}(u))\) 中的类映到 \(\delta (y^{\prime},x)\) 在 \(\mathrm{H}_{n - 1}(\mathrm{C}')\) 中的类，由此得 \(b)\) .

命题 8. --- 存在无界正合列

\[
\dots \longrightarrow \mathrm{H} _ {n} (\mathrm{C} ^ {\prime}) \xrightarrow {\mathrm{H} _ {n} (u)} \mathrm{H} _ {n} (\mathrm{C}) \xrightarrow {\mathrm{H} _ {n} (\pi)} \mathrm{H} _ {n} (\text { Con } (u)) \xrightarrow {\mathrm{H} _ {n} (\delta)} \mathrm{H} _ {n - 1} (\mathrm{C} ^ {\prime}) \longrightarrow \dots .\tag{8}
\]

事实上，根据引理3之(a)，这可由X的定理1（第30页）应用于正合列(7)得出。

推论。--- 要使u为拟同构，其充分必要条件是Con(u)具有零同调。

注：--- 考虑该图。

\[
\begin{array}{c c c c c} \dots \longrightarrow \mathrm{H} _ {n} (\mathrm{C} ^ {\prime}) & \xrightarrow {\mathrm{H} _ {n} (\tilde {u})} \mathrm{H} _ {n} (\text {Cyl} (u)) & \xrightarrow {\mathrm{H} _ {n} (\tilde {\pi})} \mathrm{H} _ {n} (\text {Con} (u)) & \xrightarrow {\partial_ {n} (\tilde {u} , \tilde {\pi})} \mathrm{H} _ {n - 1} (\mathrm{C} ^ {\prime}) \longrightarrow \dots \\ \downarrow^ {1} & \xrightarrow {\mathrm{H} _ {n} (\beta)} \downarrow & \downarrow^ {1} & \downarrow^ {1} \\ \dots \longrightarrow \mathrm{H} _ {n} (\mathrm{C} ^ {\prime}) & \xrightarrow {\mathrm{H} _ {n} (u)} \mathrm{H} _ {n} (\mathrm{C}) & \xrightarrow {\mathrm{H} _ {n} (\pi)} \mathrm{H} _ {n} (\text {Con} (u)) & \xrightarrow {\mathrm{H} _ {n} (\delta)} \mathrm{H} _ {n - 1} (\mathrm{C} ^ {\prime}) \longrightarrow \dots \end{array}
\]

其中第一行（相应地，第二行）是与正合列（6）（相应地，（7））相关联的同调正合列。映射 \(H_{n}(\beta)\) 是双射的（命题7，c）），且该图可交换，因为

\begin{enumerate}
\def\labelenumi{\alph{enumi})}
\tightlist
\item
  \(u = \beta \circ \tilde{u}\) （命题7，a）），因此 \(\mathrm{H}_n(u) = \mathrm{H}_n(\beta) \circ \mathrm{H}_n(\tilde{u})\) b） \(\mathrm{H}_{n}(\beta)=\mathrm{H}_{n}(\alpha)^{-1}\;\mathrm{and}\;\pi=\tilde{\pi}\circ\alpha\;(\mathrm{prop.~7,~a})\;\mathrm{and~c}),\;\mathrm{hence~}\mathrm{H}_{n}(\tilde{\pi})=\mathrm{H}_{n}(\pi)\circ\mathrm{H}_{n}(\beta),\) c） \(\mathrm{H}_{n}(\delta)=\partial_{n}(\tilde{u},\tilde{\pi})\) （引理3，b））。
\end{enumerate}

